\documentclass[10pt,a4paper]{amsart}
\usepackage{amsmath,amssymb,amsthm,graphicx}
% Retained minimal preamble for the Stacks Project chapter sites-cohomology.tex (Cohomology on Sites).
% Verbatim excerpt of the chapter preamble: Xy-pic setup (native preamble lines 14-18) and the native
% theorem declarations with the mathematical macros (native preamble lines 154-217). Copied mechanically;
% nothing added, removed or redefined. The wrapper class and the amsmath/amssymb/amsthm/graphicx base
% packages are supplied by the assembler, so the class/externaldocument part of the native preamble is omitted.
% For commutative diagrams we use Xy-pic
\usepackage[all]{xy}

% We use 2cell for 2-commutative diagrams.
\xyoption{2cell}
% Theorem environments.
%
\theoremstyle{plain}
\newtheorem{theorem}[subsection]{Theorem}
\newtheorem{proposition}[subsection]{Proposition}
\newtheorem{lemma}[subsection]{Lemma}

\theoremstyle{definition}
\newtheorem{definition}[subsection]{Definition}
\newtheorem{example}[subsection]{Example}
\newtheorem{exercise}[subsection]{Exercise}
\newtheorem{situation}[subsection]{Situation}

\theoremstyle{remark}
\newtheorem{remark}[subsection]{Remark}
\newtheorem{remarks}[subsection]{Remarks}

\numberwithin{equation}{subsection}

% Macros
%
\def\lim{\mathop{\mathrm{lim}}\nolimits}
\def\colim{\mathop{\mathrm{colim}}\nolimits}
\def\Spec{\mathop{\mathrm{Spec}}}
\def\Hom{\mathop{\mathrm{Hom}}\nolimits}
\def\Ext{\mathop{\mathrm{Ext}}\nolimits}
\def\SheafHom{\mathop{\mathcal{H}\!\mathit{om}}\nolimits}
\def\SheafExt{\mathop{\mathcal{E}\!\mathit{xt}}\nolimits}
\def\Sch{\mathit{Sch}}
\def\Mor{\mathop{\mathrm{Mor}}\nolimits}
\def\Ob{\mathop{\mathrm{Ob}}\nolimits}
\def\Sh{\mathop{\mathit{Sh}}\nolimits}
\def\NL{\mathop{N\!L}\nolimits}
\def\CH{\mathop{\mathrm{CH}}\nolimits}
\def\proetale{{pro\text{-}\acute{e}tale}}
\def\etale{{\acute{e}tale}}
\def\QCoh{\mathit{QCoh}}
\def\Ker{\mathop{\mathrm{Ker}}}
\def\Im{\mathop{\mathrm{Im}}}
\def\Coker{\mathop{\mathrm{Coker}}}
\def\Coim{\mathop{\mathrm{Coim}}}

% Boxtimes
%
\DeclareMathSymbol{\boxtimes}{\mathbin}{AMSa}{"02}

%
% Macros for moduli stacks/spaces
%
\def\QCohstack{\mathcal{QC}\!\mathit{oh}}
\def\Cohstack{\mathcal{C}\!\mathit{oh}}
\def\Spacesstack{\mathcal{S}\!\mathit{paces}}
\def\Quotfunctor{\mathrm{Quot}}
\def\Hilbfunctor{\mathrm{Hilb}}
\def\Curvesstack{\mathcal{C}\!\mathit{urves}}
\def\Polarizedstack{\mathcal{P}\!\mathit{olarized}}
\def\Complexesstack{\mathcal{C}\!\mathit{omplexes}}
% \Pic is the operator that assigns to X its picard group, usage \Pic(X)
% \Picardstack_{X/B} denotes the Picard stack of X over B
% \Picardfunctor_{X/B} denotes the Picard functor of X over B
\def\Pic{\mathop{\mathrm{Pic}}\nolimits}
\def\Picardstack{\mathcal{P}\!\mathit{ic}}
\def\Picardfunctor{\mathrm{Pic}}
\def\Deformationcategory{\mathcal{D}\!\mathit{ef}}
\usepackage{cleveref}
\crefname{originallemma}{original lemma}{original lemmas}
\Crefname{originallemma}{Original lemma}{Original lemmas}
\makeatletter
\newlabel{sites-modules-lemma-adjoint-pullback-pushforward-modules}{{03D7}{source}{}{sourcepointer-1}{}}
\newlabel{sites-modules-lemma-adjoint-pullback-pushforward-modules@cref}{{[originallemma][0][]03D7}{[1][0][]source}{}{}{}}
\makeatother
\title{Pushforward along a flat morphism of ringed topoi preserves K-injectivity}
\author{The Stacks Project Authors}
\date{stacks-project:sites-cohomology.tex:local-snapshot-20261011; GNU FDL 1.2 or later}
\begin{document}
\maketitle
\noindent\textit{Editorial note.} All delivered mathematical text is the author's own, copied from the pinned Stacks Project chapter file sites-cohomology.tex; no statement or proof step was rewritten, reordered or removed.\par\medskip
\noindent\textit{Editorial note.} The proof is the author's complete argument: the adjunction identification and the use of exactness of f\textasciicircum{}* are both in the counted span, and the cited adjunction is named with its chapter.\par\medskip
\noindent\textit{Editorial note.} Cross-chapter restatement: the same label occurs in chapter cohomology and in chapter more-algebra. Those statements are the special cases of this ringed-site statement, so per the lane-s4 rule only this more general occurrence is counted here.\par\medskip
\noindent\textit{Editorial note.} The single labelled reference occurring inside the counted spans points to a chapter that is not staged in this repository; it is carried in assembly.external\_references as a non-load-bearing navigation pointer whose source location is the invoking line of the pinned chapter, so no load-bearing mathematics is replaced by a stub.\par\medskip
\noindent\textit{Editorial note.} The preamble is the chapter's own retained minimal preamble (sources/stacks-sites-cohomology/retained-preamble.tex), a verbatim excerpt of the native preamble: the Xy-pic setup and the native theorem declarations with the mathematical macros.\par\medskip
\noindent\textit{Allowed standard prerequisites (editorial).}\begin{itemize}
\item Flat morphism of ringed topoi. Same chapter, the definition of a flat morphism of ringed topoi, i.e. exactness of the pullback f\textasciicircum{}*. Supplies the hypothesis that f\textasciicircum{}* is exact, which is what turns the adjunction isomorphism into a K-injectivity statement.
\item K-injective complex of modules on a ringed topos. Same chapter, the defining property of K-injectivity in terms of homs in the homotopy category. Supplies the property to be verified for f\_* I\textasciicircum{}bullet, tested against homotopy classes of maps.
\item Adjunction between pullback and pushforward of modules on ringed topoi, at the level of complexes. Modules on Sites, the adjunction Hom\_\{K(O\_D)\}(F\textasciicircum{}bullet, f\_* I\textasciicircum{}bullet) = Hom\_\{K(O\_C)\}(f\textasciicircum{}* F\textasciicircum{}bullet, I\textasciicircum{}bullet). Supplies the identification of hom groups with which the counted proof concludes.
\end{itemize}\medskip
\ifdefined\hypertarget\hypertarget{sourcepointer-1}{}\fi
\noindent\textit{Source reference.} Original lemma 03D7: sources/stacks-sites-cohomology/pinned-original-source.tex, source lines 4611--4611. Cross-chapter navigation pointer invoked at proof lines 4609-4612, where the author uses the adjunction between f\textasciicircum{}* and f\_* on complexes of modules. The target chapter (sites.tex, Modules on Sites) is not staged in this repository, so the pointer is registered with its source location on the invoking line of the pinned chapter; it only resolves the reference for the reader and reproduces no mathematics, and no bibliography citation is involved.\par\medskip
\begin{lemma}
\label{lemma-K-injective-flat}
Let $f : (\Sh(\mathcal{C}), \mathcal{O}_\mathcal{C}) \to
(\Sh(\mathcal{D}), \mathcal{O}_\mathcal{D})$ be a flat morphism
of ringed topoi. If $\mathcal{I}^\bullet$ is a K-injective
complex of $\mathcal{O}_\mathcal{C}$-modules, then
$f_*\mathcal{I}^\bullet$ is K-injective
as a complex of $\mathcal{O}_\mathcal{D}$-modules.
\end{lemma}

\begin{proof}
This is true because
$$
\Hom_{K(\mathcal{O}_\mathcal{D})}(\mathcal{F}^\bullet, f_*\mathcal{I}^\bullet)
=
\Hom_{K(\mathcal{O}_\mathcal{C})}(f^*\mathcal{F}^\bullet, \mathcal{I}^\bullet)
$$
by
Modules on Sites, Lemma
\ref{sites-modules-lemma-adjoint-pullback-pushforward-modules}
and the fact that $f^*$ is exact as $f$ is assumed to be flat.
\end{proof}

\end{document}
